
{\sffamily

\thispagestyle{empty}

\begin{flushleft}
\Large{\textbf{CZECH TECHNICAL UNIVERSITY IN PRAGUE}}\\
\Large{\textbf{FACULTY OF ELECTRICAL ENGINEERING}}\\
\Large{\textbf{DEPARTMENT OF COMPUTER SCIENCE}}
\end{flushleft}

\vspace*{1em}

\begin{flushleft}
\Large{\textbf{Algorithms and Iterative Methods for Infinite Games}}\\
\Large{\textbf{Doctoral Thesis}}
\end{flushleft}

\vspace*{1em}

\begin{flushleft}
\Large{\textbf{Tomáš Votroubek}}
\end{flushleft}

\vspace*{1em}

\begin{flushleft}
Prague, February, 2026\\
\end{flushleft}

\vspace*{1em}

\begin{flushleft}
\textbf{DOCTORAL STUDY PROGRAM: Informatics}\\
\textbf{BRANCH OF STUDY: DCINFCSC -- Informatics}
\end{flushleft}

\vspace*{1em}

\begin{flushleft}
\textbf{SUPERVISOR:  Kroupa Tomáš doc. Ing. Ph.D.}\\
\textbf{SUPERVISOR-SPECIALIST: ---}
\end{flushleft}

\pagebreak
\thispagestyle{empty}

{\small
\paragraph{Abstract}
In practice, the inherent complexity of game-theoretical concepts often necessitates simplifying assumptions to model interactions between agents. When it is possible to map our problem to one of the well-studied classes of games, established solution concepts can usually be practically computed. Interesting interactions in the real world care little about game theory, however, and they appear arbitrarily far from any of the convenient classes of games. In scenarios where the reductions introduced by such mapping would be too drastic, we may still gain some insight into players' behavior with ad hoc heuristics, such as by discretizing the action spaces. While these approaches can be useful, they may lack robustness and guarantees.

In this thesis, we explore methods capable of providing guarantees or certificates on the quality of their solutions across a broad range of games. Consequently, we explore methods for global optimization, which appears as a very common sub-task in game theory. The~ability to check global optimality is, for example, useful in answering whether an agent in an interaction is behaving rationally.

Finding global optima of functions over a set in the absence of any assumptions is a hard problem; even approximating such solutions is hard. Nevertheless, there exist classes of problems where global optimization is feasible. We demonstrate how Lasserre's hierarchies and the Spatial Branch and Bound method can be effectively applied to solve tasks in game theory and robotics. Our research indicates that the practicality of global methods remains largely unexplored despite the potential.

\paragraph{Keywords:} Game Theory, Infinite Games, $\epsilon$-Nash Equilibria, Inverse Kinematics, Global Optimality, Lasserre Hierarchies, Branch and Bound

\paragraph{Abstrakt}
Z praktických důvodů nejsou interakce mezi agenty vždy modelovány zcela věrně realitě. Pokud lze daný problém formulovat jako některou ze známých tříd her, nemusí být obtížné najít jeho řešení. Interakce v reálném světě ale právě do těch snadno řešitelných tříd často nezapadají. I v takových případech jsme schopni hledat alespoň heuristická řešení, například diskretizací strategických prostorů hráčů. S každou heuristikou a každým předpokladem přidaným nad rámec původního problému ale přicházíme o záruky na kvalitu řešení.

V této práci zkoumáme metody schopné spočítat řešení se zaručenou kvalitou pro širokou škálu her. Nevyhnutelně se proto zabýváme také metodami globální optimalizace, která představuje velmi běžný dílčí úkol v teorii her. Schopnost ověřit globální optimalitu je užitečná zejména pokud chceme vědět, zda se agent v dané situaci chová racionálně.

Nalezení globálního optima funkcí nad množinou bez jakýchkoli předpokladů je náročný problém; a i aproximace takových řešení je náročná. Nicméně existují třídy problémů, u nichž je globální řešení nalézt lze. V této práci ukážeme příklady použití Lasserreových hierarchií a metody Spatial Branch and Bound pro efektivní řešení úloh v teorii her a robotice. Náš výzkum naznačuje, že praktičnost globálních metod zůstává navzdory jejich potenciálu z velké části neprozkoumaná.

\paragraph{Klíčová slova:} Teorie her, Nekonečné hry, $\epsilon$-Nashovy rovnováhy, inverzní kinematika, globální optimalita, Lasserreovy hierarchie, metoda větví a mezí
}

\pagebreak

\thispagestyle{empty}

\paragraph{Publications}
\begin{itemize}
    \item \fullcite{cvutIK24}
    \item \fullcite{cvutMO23}
    \item \fullcite{cvutBool}
    \item \fullcite{cvutNet21}
\end{itemize}

}

